% Supplementary Table: absolute fit-quality metrics for the ten
% lowest-objective separate in vivo fits.
%
% PURPOSE
% This table documents the per-fit values underlying the first Results
% paragraph of "In vivo fits capture cohort-scale behavior while retaining
% multiple compatible parameterizations." It expands every reported range
% into the ten selected numerical fits.
%
% PRIMARY RESULT ROOT
%   /Users/4482173/Documents/GitHub/soft_coupling/oxygen/results/
%   fit_invivo_O2_buffering_500seed/
%
% PER-SEED SOURCE FILES
% 1. Objective and optimizer summary:
%    seed*/fit_summary.tsv
% 2. Longitudinal tumor-burden observations and predictions:
%    seed*/burden_fit.tsv
% 3. Terminal chromosome-number distributions:
%    seed*/terminal_ploidy_fit.tsv
% 4. Observed terminal necrosis fractions and matching harvest/day keys:
%    seed*/necrosis_fit.tsv
%
% TOP-TEN SELECTION AND PORTABLE AUDIT FILES
% 1. Ranked source-directory index:
%    /Users/4482173/Documents/GitHub/HypoxiaLTEEFigures/revised/SuppFiles/
%    round1/miningcloneid_review_v2/analysis_tables/
%    source_top10_index_localized.csv
%    Selection: result_group = fit_invivo_O2_buffering_500seed and
%    rank = 1,...,10 after ascending objective ranking.
% 2. Exact per-seed values used in this table:
%    /Users/4482173/Documents/GitHub/HypoxiaLTEEFigures/revised/SuppFiles/
%    round1/miningcloneid_review_v2/analysis/
%    invivo_top10_fit_quality.csv
% 3. Extended terminal-distribution cross-check:
%    /Users/4482173/Documents/GitHub/HypoxiaLTEEFigures/revised/SuppFiles/
%    round1/miningcloneid_review_v2/analysis_tables/
%    invivo_terminal_distribution_fit_quality_extended.csv
% 4. Recalculation script:
%    /Users/4482173/Documents/GitHub/HypoxiaLTEEFigures/revised/SuppFiles/
%    round1/miningcloneid_review_v2/scripts/reanalyse_results.py
%    Relevant block: "# ---------------- in vivo ----------------".
%
% OBJECTIVE
% The displayed objective is objective_total from fit_summary.tsv and agrees
% with the ranked objective in source_top10_index_localized.csv.
%
% POSITIVE-OBSERVATION LOG-BURDEN RMSE
% From burden_fit.tsv, retain rows satisfying:
%   day > 0 and obs_burden > 0.
% For each retained observation i, define:
%   r_i = pred_log_burden_i - obs_log_burden_i.
% Then:
%   burden log-RMSE = sqrt(mean_i(r_i^2)).
%
% TUMOR-BALANCED LOG-BURDEN RMSE
% Within each harvest/tumor h, first calculate mean_{i in h}(r_i^2).
% Give every harvest equal weight, then calculate:
%   tumor-balanced log-RMSE
%     = sqrt(mean_h[mean_{i in h}(r_i^2)]).
% This differs from the observation-level RMSE because tumors contribute
% equally rather than in proportion to their number of burden measurements.
%
% TERMINAL CHROMOSOME-NUMBER METRICS
% For each harvest in terminal_ploidy_fit.tsv:
% 1. Normalize pred_fraction and obs_count separately to sum to one.
% 2. Predicted mean N = sum_N(N * normalized pred_fraction_N).
% 3. Observed mean N = sum_N(N * normalized obs_count_N).
% 4. Mean-N absolute error = abs(predicted mean N - observed mean N).
% 5. Total-variation distance
%      = 0.5 * sum_N(abs(pred_fraction_N - obs_fraction_N)).
% 6. Wasserstein-1 distance on the consecutive one-chromosome N grid
%      = sum_N(abs(cumsum(pred_fraction)_N - cumsum(obs_fraction)_N)).
% The table reports the arithmetic mean of each metric across harvests.
% Mean-N MAE and Wasserstein-1 are in chromosome-number units; total
% variation is dimensionless.
%
% RECONSTRUCTED TERMINAL NECROSIS-FRACTION MAE
% The stored pred_necrosis_fraction column in necrosis_fit.tsv is missing for
% all rows in these exports. Therefore, for each necrosis observation, match
% burden_fit.tsv by harvest and day and reconstruct:
%   predicted necrosis fraction
%     = pred_burden_dead_total_volume_mm3 / pred_burden_volume_mm3.
% Then calculate:
%   necrosis-fraction MAE
%     = mean(abs(predicted fraction - observed necrosis fraction)).
% This reconstruction is why the Results and table use the word
% "reconstructed." It is an audit of stored model states, not a repair of the
% missing row-level prediction column in necrosis_fit.tsv.
%
% VERIFIED UNROUNDED RANGES
% Objective: 14.1193277940156--14.1747852749872.
% Positive-observation log-burden RMSE:
%   0.662868705233521--0.679850558662106.
% Tumor-balanced log-burden RMSE:
%   0.628359037957605--0.643940087623863.
% Terminal mean-N MAE:
%   2.40731057368453--3.71221228124632 chromosomes.
% Mean terminal Wasserstein-1:
%   4.18402106806819--5.23917736605831 chromosomes.
% Mean terminal total variation:
%   0.681084475678739--0.713875341871515.
% Reconstructed terminal necrosis-fraction MAE:
%   0.0644328419561802--0.102546103165418.
%
% INTERPRETIVE LIMIT
% These rows are optimizer-selected numerical fits, not biological
% replicates, posterior samples, or a confidence interval. The metrics
% quantify in-sample descriptive fit quality under the fitted model.

\begin{table}[!htbp]
\centering
\scriptsize
\setlength{\tabcolsep}{3.5pt}
\caption{Absolute fit-quality metrics for the ten lowest-objective separate \textit{in vivo} fits. Burden errors are calculated on the log-burden scale. Terminal chromosome-number metrics and reconstructed necrosis-fraction MAE are averaged equally across tumors/harvests.}
\label{tab:invivo_top10_fit_quality}
\resizebox{\textwidth}{!}{%
\begin{tabular}{@{}rlrrrrrrr@{}}
\toprule
Rank & Fit & Objective &
\shortstack{Burden\\log-RMSE} &
\shortstack{Tumor-balanced\\log-RMSE} &
\shortstack{Terminal mean-\(N\)\\MAE} &
\shortstack{Terminal\\Wasserstein-1} &
\shortstack{Terminal total\\variation} &
\shortstack{Reconstructed necrosis-\\fraction MAE} \\
\midrule
1  & \texttt{seed25}  & 14.1193 & 0.6660 & 0.6317 & 2.41 & 4.51 & 0.697 & 0.091 \\
2  & \texttt{seed366} & 14.1340 & 0.6680 & 0.6341 & 2.54 & 4.75 & 0.708 & 0.078 \\
3  & \texttt{seed292} & 14.1372 & 0.6724 & 0.6384 & 2.79 & 4.59 & 0.698 & 0.064 \\
4  & \texttt{seed392} & 14.1406 & 0.6784 & 0.6426 & 3.33 & 4.92 & 0.689 & 0.093 \\
5  & \texttt{seed90}  & 14.1524 & 0.6799 & 0.6439 & 3.65 & 4.86 & 0.684 & 0.092 \\
6  & \texttt{seed391} & 14.1553 & 0.6797 & 0.6439 & 3.38 & 4.75 & 0.682 & 0.091 \\
7  & \texttt{seed264} & 14.1558 & 0.6735 & 0.6400 & 2.56 & 4.18 & 0.682 & 0.079 \\
8  & \texttt{seed109} & 14.1724 & 0.6629 & 0.6284 & 2.85 & 4.30 & 0.681 & 0.103 \\
9  & \texttt{seed322} & 14.1724 & 0.6677 & 0.6331 & 2.52 & 4.86 & 0.714 & 0.078 \\
10 & \texttt{seed26}  & 14.1748 & 0.6784 & 0.6425 & 3.71 & 5.24 & 0.698 & 0.094 \\
\bottomrule
\end{tabular}%
}
\par\vspace{0.5em}
\begin{minipage}{0.96\textwidth}
\footnotesize
\textit{Note:} Positive-observation burden RMSE weights every retained measurement equally; tumor-balanced RMSE first averages squared residuals within each tumor and then weights tumors equally. Wasserstein-1 and mean-\(N\) MAE are reported in chromosomes. Necrosis predictions were reconstructed from matched dead-total and total predicted burden volumes because the exported row-level predicted-necrosis column was missing. The ten optimizer-selected fits are not biological replicates or confidence-interval samples.
\end{minipage}
\end{table}
